\scanpage{75}

% L01 -- top
Posant

% L02 -- top
\[
\omega_V=a_0\wedge\cdots\wedge a_n
\]

% L03 -- top
on a

% L04 -- upper third, discriminant
\[
\delta(f_V)=\underbrace{\beta_0^n\beta_1^{n-1}\cdots\beta_{n-1}}_{B}
 \rho\,\omega_V^{\otimes-2}=B\rho\,\omega_V^{\otimes-2}
\]

% L05 -- below discriminant
($f$ lisse sss $B$ inv. i.e. les $\beta_i$ inv \quad $0\le i\le1$.)

% L06 -- upper middle
Considérons alors

% L07 -- upper middle, parity cases
\[
\lambda=f_V(\widetilde{s}_0)=\deleted{\unsure{-}}+
\frac{\sigma'\rho'}{\rho^2}
=\left\{\begin{aligned}
&\left(\begin{gathered}n\text{ pair \unsure{avec|donc} }\rho=2\rho'\\(2\text{ inv!})\end{gathered}\right)
&&+\frac{\sigma'}{4\rho'}\\
&n\text{ impair }\ \rho=\rho'&&+\frac{\sigma'}{\rho'}
\end{aligned}\right.
\]

% L08 -- middle, quadric
L'\deleted{\unsure{a}} \deleted{\unsure{hyp}}hyperquadrique affine circonscrite au polyèdre régulier est alors donnée par

% L09 -- middle
\[
f_V-\lambda=0
\]

% L10 -- middle, smoothness
Pour qu'elle soit \deleted{affine} \inserted{\unsure{projectivement|respectivement}} lisse il f. et s. que les $\beta_i$ et $\sigma'$ soient inv (c'est donc le cas \uline{\unsure{birégulier|singulier}})

% L11 -- lower third, D
(D) \uline{NB} Supposons donnés $V$, $f_V$ et $\varphi_G^V$, représentation \deleted{\unsure{ortho.|orth.}}

% L12 -- lower third, action
\[
G\longrightarrow \mathrm O(f_{\unsure{V|W}})
\]

% L13 -- below action
$1^\circ$ transformant les $\tau_i$ en réflexions \uline{orth.} et satisfaisant de plus $2^\circ)$ $3^\circ)$ plus haut.

% L14 -- lower third, revised wording
Ceci \uline{implique} \deleted{\unsure{satisfaisant de plus}},

% L15 -- above foot, replacement condition
\[
\deleted{\unsure{4^\circ)}}\quad
\bigl[\deleted{f\mid(D_0=\bigcap_{1\le i\le n}H_i)}
\inserted{f\mid L_0}\ne0\quad\text{en \unsure{tte|ts} fibre.}
\]

% L16 -- lower left, diagonal, double-headed pointer to condition
% slanted
\margin{$\longleftrightarrow$ (i.e. $\tau_0$ est une réflexion \unsure{str.} orth. de $V$) \deleted{\unsure{Ceci} Cette condition est-elle une \unsure{conséq.} des autres ?}}

% L17 -- foot, continuation
Alors \unsure{ça|$s_0$} provient \unsure{loc|doc|donc} fppf d'une situation comme plus haut. Il suffit de choisir une section \inserted{$a_0$} de $D_0$ telle que $f_V(a_0)=1$
% archive sheet number: 74
